\section{The two exact sequences}\label{sec:exact}
% ===============================================================

Two questions organise what follows.
\begin{itemize}[leftmargin=1.5em,itemsep=0.15em]
\item Which divisors come from meromorphic functions --- and how much of a
function is remembered by its divisor? This is answered in this section, by the
first exact sequence, and completed by Theorem~\ref{thmA}.
\item How far is an arbitrary divisor from being principal? The second exact
sequence sets up the quotient that measures this; the measurement itself is
carried out in Section~\ref{sec:cl}.
\end{itemize}

\subsection{The kernel of the divisor map}

\begin{smjlemma}\label{lem:constant}
Every function holomorphic on all of $\PP$ is constant.
\end{smjlemma}

\begin{proof}
Let $h$ be holomorphic on $\PP$. Then $|h|$ is continuous on the compact space
$\PP$ (Proposition~\ref{prop:compact}), so it attains a maximum at some
$p_0 \in \PP$. Choose a chart around $p_0$; in it $h$ is a holomorphic function
of one variable whose modulus has an interior maximum, so by the maximum modulus
principle \cite[Ch.~4]{Ahlfors} $h$ is constant, say equal to $c$, on a
neighbourhood of $p_0$. The set $\{p \in \PP : h(p) = c\}$ is therefore nonempty
and open; it is closed by continuity; and $\PP$ is connected
(Proposition~\ref{prop:compact}). Hence it is all of $\PP$.
\end{proof}

\begin{smjproposition}\label{prop:kernel}
$\ker(\divv) = \CC^\times$.
\end{smjproposition}

\begin{proof}
If $c \in \CC^\times$ then $\divv(c) = 0$ by Proposition~\ref{prop:divconst}, so
$\CC^\times \subseteq \ker(\divv)$. Conversely let $f \in \CC(z)^\times$ with
$\divv(f) = 0$. Then $\ord_p(f) = 0$ for every $p \in \PP$, so $f$ has no poles
and is holomorphic on all of $\PP$. By Lemma~\ref{lem:constant} $f$ is constant,
and it is nonzero, so $f \in \CC^\times$.
\end{proof}

\begin{remark}
Equivalently: two rational functions have the same divisor if and only if they
differ by a nonzero constant multiple. Indeed $\divv(f) = \divv(g)$ gives
$\divv(f/g) = 0$, so $f/g \in \CC^\times$. A divisor therefore remembers a
function exactly up to scale --- no more and no less.
\end{remark}

\subsection{The first exact sequence}

\begin{smjproposition}\label{prop:seq1}
The sequence
\[
1 \longrightarrow \CC^\times \longrightarrow \CC(z)^\times
\stackrel{\divv}{\longrightarrow} \Prin(\PP) \longrightarrow 1
\]
is exact.
\end{smjproposition}

\begin{proof}
The first map is the inclusion of constants, which is injective. Exactness at
$\CC(z)^\times$ is Proposition~\ref{prop:kernel}. Exactness at $\Prin(\PP)$ is
the surjectivity of $\divv$ onto its image, which holds by
Definition~\ref{def:prin}.
\end{proof}

\begin{remark}
By the First Isomorphism Theorem this says
$\Prin(\PP) \cong \CC(z)^\times / \CC^\times$: the group of principal divisors
\emph{is} the group of nonzero rational functions modulo scaling.
\end{remark}

\subsection{Linear equivalence and the divisor class group}

\begin{definition}\label{def:lineq}
Divisors $D, E \in \Div(\PP)$ are \emph{linearly equivalent}, written
$D \sim E$, if $D - E \in \Prin(\PP)$.
\end{definition}

\begin{smjproposition}
Linear equivalence is an equivalence relation on $\Div(\PP)$.
\end{smjproposition}

\begin{proof}
This is exactly the statement that $\Prin(\PP)$ is a subgroup of the abelian
group $\Div(\PP)$ (Definition~\ref{def:prin}). Reflexivity: $D - D = 0 =
\divv(1)$. Symmetry: if $D - E = \divv(f)$ then
$E - D = -\divv(f) = \divv(f^{-1})$. Transitivity: if $D - E = \divv(f)$ and
$E - F = \divv(g)$ then $D - F = \divv(f) + \divv(g) = \divv(fg)$, using
Proposition~\ref{prop:divhom} throughout.
\end{proof}

\begin{definition}
The \emph{divisor class group} of $\PP$ is the quotient
\[
\Cl(\PP) := \Div(\PP)/\Prin(\PP),
\]
whose elements are the linear equivalence classes $[D]$.
\end{definition}

\subsection{The second exact sequence}

\begin{smjproposition}
The sequence
$0 \to \Prin(\PP) \to \Div(\PP) \to \Cl(\PP) \to 0$ is exact.
\end{smjproposition}

\begin{proof}
The first map is the inclusion of a subgroup, hence injective. The second is the
quotient map $q : D \mapsto [D]$, which is surjective by construction and has
kernel $\Prin(\PP)$ by the definition of a quotient group.
\end{proof}

% ===============================================================
